\section{敏感性分析}

\subsection{参数敏感性}

对关键参数在基准值附近 $\pm 20\%$ 范围内进行扰动测试。各参数敏感性均较低（$R^2$ 最大变化仅 1.5\%），模型具有较强的鲁棒性。

\subsection{关键假设扰动与对照实验}

% 对动态 LTM 中每条带【关键】标签的假设（放置标签【全文】/【问题N】后追加
% 【关键】），给出扰动或对照实验并引用结论，回答「该假设若不成立，结论是否
% 依然成立」。
% 【全文】【关键】假设的实验写在本章；【问题N】【关键】假设的实验可写在该问题
% 章节，并在此引用其结论。
% 例：对假设 H（...）分别做 ±20% 扰动与「去除该假设」对照，结果变化不超过 X%，
% 关键结论保持成立，说明结论对假设 H 稳健。
